% =====================================================================
%  preamble.tex — 导言区
%  《三价图的拉格朗日 Floer 理论与曲线的同调镜像对称》中文译注本
%  底本：D. Auroux, A. I. Efimov, L. Katzarkov,
%        arXiv:2107.01981v3（Sel. Math. New Ser. 30 (2024), Paper 95）
%  编译：XeLaTeX，两遍；本文件由主译维护。
% =====================================================================
\documentclass[11pt,a4paper]{ctexart}

\usepackage[a4paper,margin=2.6cm,top=2.9cm,bottom=2.7cm]{geometry}
\usepackage{amsmath,amssymb,amsthm,mathtools}
\usepackage{mathrsfs}
\usepackage{graphicx}
\usepackage{array,booktabs,longtable,multirow}
\usepackage{caption}
\usepackage{enumitem}
\usepackage{xcolor}
\usepackage{fancyhdr}
\usepackage{tikz}
\usetikzlibrary{arrows,backgrounds,fit,positioning,shapes.symbols,chains}
\usepackage[all,cmtip,matrix,arrow]{xy}
\xyoption{rotate}
\xyoption{pdf}

% ---------------- 定理环境（与原文相同：按节编号，共享计数） ----------------
\newtheorem{theo}{定理}[section]
\newtheorem{lemma}[theo]{引理}
\newtheorem{defi}[theo]{定义}
\newtheorem{prop}[theo]{命题}
\newtheorem{cond}[theo]{条件}
\newtheorem{conj}[theo]{猜想}
\newtheorem{cor}[theo]{推论}
\newtheorem{remark}[theo]{注}
\newtheorem{property}[theo]{性质}
\newtheorem{example}[theo]{例}
\numberwithin{equation}{section}

% ---------------- 宏（与原文一致，块文件只许用这些） ----------------
\mathchardef\mhyphen="2D
\newcommand{\CP}{\mathbb{CP}}
\newcommand{\A}{\mathbb{A}}
\newcommand{\bbI}{\mathbb{I}}
\newcommand{\N}{\mathbb{N}}
\newcommand{\R}{\mathbb{R}}
\newcommand{\C}{\mathbb{C}}
\newcommand{\Z}{\mathbb{Z}}
\newcommand{\Q}{\mathbb{Q}}
\newcommand{\bG}{\mathbb{G}}
\newcommand{\bL}{\mathbb{L}}
\newcommand{\bH}{\mathbf{H}}
\newcommand{\bC}{\mathbf{C}}
\newcommand{\bD}{\mathbf{D}}
\newcommand{\bR}{\mathbf{R}}
\newcommand{\bY}{\mathbf{Y}}
\newcommand{\PP}{\mathbb{P}}
\newcommand{\FF}{\mathbb{F}}
\newcommand{\B}{\mathcal{B}}
\newcommand{\E}{\mathcal{E}}
\newcommand{\F}{\mathcal{F}}
\newcommand{\G}{\mathcal{G}}
\providecommand{\O}{}\renewcommand{\O}{\mathcal{O}}
\newcommand{\Pc}{\mathcal{P}}
\providecommand{\D}{}\renewcommand{\D}{\mathcal{D}}
\newcommand{\W}{\mathcal{W}}
\newcommand{\cE}{\mathcal{E}}
\newcommand{\cX}{\mathcal{X}}
\newcommand{\cM}{\mathcal{M}}
\newcommand{\cN}{\mathcal{N}}
\newcommand{\cK}{\mathcal{K}}
\newcommand{\cL}{\mathcal{L}}
\newcommand{\cU}{\mathcal{U}}
\newcommand{\cW}{\mathcal{W}}
\newcommand{\cD}{\mathcal{D}}
\newcommand{\cB}{\mathcal{B}}
\newcommand{\cP}{\mathcal{P}}
\newcommand{\mX}{\mathfrak{X}}
\newcommand{\cC}{\mathcal{C}}
\newcommand{\cF}{\mathcal{F}}
\newcommand{\cO}{\mathcal{O}}
\newcommand{\id}{\operatorname{id}}
\newcommand{\val}{\operatorname{val}}
\newcommand{\im}{\operatorname{im}}
\DeclareMathOperator{\supp}{supp}
\DeclareMathOperator{\pr}{pr}
\DeclareMathOperator{\Spec}{Spec}
\DeclareMathOperator{\Spf}{Spf}
\newcommand{\can}{\operatorname{can}}
\newcommand{\mk}{\mathbf{k}}
\newcommand{\Sch}{\mathbf{Sch}}
\newcommand{\m}{\mathfrak{m}}
\newcommand{\hto}{\hookrightarrow}
\newcommand{\xto}{\xrightarrow}
\newcommand{\prarr}{\rightrightarrows}
\newcommand{\Tw}{\operatorname{Tw}}
\newcommand{\crit}{\operatorname{crit}}
\newcommand{\Coh}{\operatorname{Coh}}
\newcommand{\Perf}{\operatorname{Perf}}
\newcommand{\Ext}{\operatorname{Ext}}
\newcommand{\Hom}{\operatorname{Hom}}
\newcommand{\End}{\operatorname{End}}
\newcommand{\Aut}{\operatorname{Aut}}
\newcommand{\NS}{\operatorname{NS}}
\newcommand{\Pic}{\operatorname{Pic}}
\newcommand{\rad}{\operatorname{rad}}
\newcommand{\gr}{\operatorname{gr}}
\newcommand{\hd}{\operatorname{hd}}
\newcommand{\db}[1]{\bD^b(#1)}
\newcommand{\coh}{\operatorname{coh}}
\newcommand{\lto}{\longrightarrow}
\newcommand{\wt}{\widetilde}
\newcommand{\pre}{\operatorname{pre-tr}}
\providecommand{\h}{}\renewcommand{\h}{\operatorname{h}}
\DeclareMathOperator*{\colim}{colim}
\DeclareMathOperator*{\hocolim}{hocolim}
\DeclareMathOperator*{\holim}{holim}
\newcommand{\tens}[1]{\mathbin{\mathop{\otimes}\displaylimits_{#1}}}
\newcommand{\Ltens}[1]{\mathbin{\mathop{\otimes}\displaylimits^{\bL}_{#1}}}

% ---------------- 版式 ----------------
\setlength{\parindent}{2em}
\linespread{1.28}
\setlist{nosep,leftmargin=2.2em}
\renewcommand{\arraystretch}{1.15}

\pagestyle{fancy}
\fancyhf{}
\fancyhead[L]{\small\kaishu 三价图的拉格朗日 Floer 理论与曲线的同调镜像对称}
\fancyhead[R]{\small\thepage}
\renewcommand{\headrulewidth}{0.4pt}
\fancypagestyle{plain}{\fancyhf{}\fancyfoot[C]{\small\thepage}\renewcommand{\headrulewidth}{0pt}}

\ctexset{
  section/format = {\Large\bfseries\heiti},
  subsection/format = {\large\bfseries\heiti},
  subsubsection/format = {\normalsize\bfseries\heiti},
}

\usepackage[colorlinks=true,linkcolor=blue!45!black,citecolor=blue!45!black,
            urlcolor=blue!45!black,
            bookmarksnumbered=true,bookmarksopen=true,
            pdftitle={三价图的拉格朗日 Floer 理论与曲线的同调镜像对称（中文译注本）},
            pdfauthor={D. Auroux, A. I. Efimov, L. Katzarkov；译者：数译}]{hyperref}
